
\documentclass{article}
\usepackage{colortbl}
\usepackage{makecell}
\usepackage{multirow}
\usepackage{supertabular}

\begin{document}

\newcounter{utterance}

\twocolumn

{ \footnotesize  \setcounter{utterance}{1}
\setlength{\tabcolsep}{0pt}
\begin{supertabular}{c@{$\;$}|p{.15\linewidth}@{}p{.15\linewidth}p{.15\linewidth}p{.15\linewidth}p{.15\linewidth}p{.15\linewidth}}

    \# & $\;$A & \multicolumn{4}{c}{Game Master} & $\;\:$B\\
    \hline 

    \theutterance \stepcounter{utterance}  

    & & \multicolumn{4}{p{0.6\linewidth}}{\cellcolor[rgb]{0.9,0.9,0.9}{%
	\makecell[{{p{\linewidth}}}]{% 
	  \tt {\tiny [A$\langle$GM]}  
	 Du nimmst an einem kollaborativen Verhandlungspiel Teil.\\ \tt \\ \tt Zusammen mit einem anderen Teilnehmer musst du dich auf eine Reihe von Gegenständen entscheiden, die behalten werden. Jeder von euch hat eine persönliche Verteilung über die Wichtigkeit der einzelnen Gegenstände. Jeder von euch hat eine eigene Meinung darüber, wie wichtig jeder einzelne Gegenstand ist (Gegenstandswichtigkeit). Du kennst die Wichtigkeitsverteilung des anderen Spielers nicht. Zusätzlich siehst du, wie viel Aufwand jeder Gegenstand verursacht.  \\ \tt Ihr dürft euch nur auf eine Reihe von Gegenständen einigen, wenn der Gesamtaufwand der ausgewählten Gegenstände den Maximalaufwand nicht überschreitet:\\ \tt \\ \tt Maximalaufwand = 8145\\ \tt \\ \tt Hier sind die einzelnen Aufwände der Gegenstände:\\ \tt \\ \tt Aufwand der Gegenstände = {"B54": 963, "C72": 204, "B49": 418, "C99": 961, "A69": 238, "B36": 623, "A23": 44, "C77": 633, "C61": 247, "B85": 645, "A90": 230, "C23": 250, "C22": 731, "C51": 405, "B27": 389, "A42": 216, "C13": 637, "C12": 156, "A65": 738, "B07": 307, "A48": 762, "A59": 737, "C98": 901, "A84": 369, "C33": 2, "A87": 730, "C90": 722, "A74": 704, "A50": 315, "B44": 455, "A99": 510, "A34": 175, "B17": 692, "A62": 150, "C60": 32}\\ \tt \\ \tt Hier ist deine persönliche Verteilung der Wichtigkeit der einzelnen Gegenstände:\\ \tt \\ \tt Werte der Gegenstandswichtigkeit = {"B54": 1059, "C72": 300, "B49": 514, "C99": 1057, "A69": 334, "B36": 719, "A23": 140, "C77": 729, "C61": 343, "B85": 741, "A90": 326, "C23": 346, "C22": 827, "C51": 501, "B27": 485, "A42": 312, "C13": 733, "C12": 252, "A65": 834, "B07": 403, "A48": 858, "A59": 833, "C98": 997, "A84": 465, "C33": 98, "A87": 826, "C90": 818, "A74": 800, "A50": 411, "B44": 551, "A99": 606, "A34": 271, "B17": 788, "A62": 246, "C60": 128}\\ \tt \\ \tt Ziel:\\ \tt \\ \tt Dein Ziel ist es, eine Reihe von Gegenständen auszuhandeln, die dir möglichst viel bringt (d. h. Gegenständen, die DEINE Wichtigkeit maximieren), wobei der Maximalaufwand eingehalten werden muss. Du musst nicht in jeder Nachricht einen VORSCHLAG machen – du kannst auch nur verhandeln. Alle Taktiken sind erlaubt!\\ \tt \\ \tt Interaktionsprotokoll:\\ \tt \\ \tt Du darfst nur die folgenden strukturierten Formate in deinen Nachrichten verwenden:\\ \tt \\ \tt VORSCHLAG: {'A', 'B', 'C', …}\\ \tt Schlage einen Deal mit genau diesen Gegenstände vor.\\ \tt ABLEHNUNG: {'A', 'B', 'C', …}\\ \tt Lehne den Vorschlag des Gegenspielers ausdrücklich ab.\\ \tt ARGUMENT: {'...'}\\ \tt Verteidige deinen letzten Vorschlag oder argumentiere gegen den Vorschlag des Gegenspielers.\\ \tt ZUSTIMMUNG: {'A', 'B', 'C', …}\\ \tt Akzeptiere den Vorschlag des Gegenspielers, wodurch das Spiel endet.\\ \tt STRATEGISCHE ÜBERLEGUNGEN: {'...'}\\ \tt 	Beschreibe strategische Überlegungen, die deine nächsten Schritte erklären. Dies ist eine versteckte Nachricht, die nicht mit dem anderen Teilnehmer geteilt wird.\\ \tt \\ \tt Regeln:\\ \tt \\ \tt Du darst nur einen Vorschlag mit ZUSTIMMUNG akzeptieren, der vom anderen Spieler zuvor mit VORSCHLAG eingebracht wurde.\\ \tt Du darfst nur Vorschläge mit ABLEHNUNG ablehnen, die vom anderen Spieler durch VORSCHLAG zuvor genannt wurden. \\ \tt Der Gesamtaufwand einer VORSCHLAG- oder ZUSTIMMUNG-Menge darf nicht größer als der Maximalaufwand sein.  \\ \tt Offenbare deine versteckte Wichtigkeitsverteilung nicht.\\ \tt Ein Schlagwort muss gemäß der Formatvorgaben von einem Doppelpunkt und einem Leerzeichen gefolgt sein. Das Argument muss eine Python-Menge mit 0 oder mehr Strings sein.  \\ \tt Die Form muss also SCHLAGWORT: {...} sein.\\ \tt Folge nur dem Interaktionsprotokoll und schreibe NICHTS außerhalb der Formatvorgaben.\\ \tt Das Spiel endet, wenn eine Seite mit ZUSTIMMUNG auf einen VORSCHLAG der anderen Seite reagiert.  \\ \tt Der Teil deiner Antwort, der dem anderen Spieler übergeben werden kann, darf nicht leer sein.  \\ \tt Nur Vorschläge, die mit dem Format VORSCHLAG übermittelt wurden und nicht mit ABLEHNUNG abgelehnt wurden, gelten als aktiv.  \\ \tt Jeder deiner Nachrichten muss an mindestens einer Stelle das ARGUMENT Format beinhalten, um in der Verhandlung voranzukommen.\\ \tt Du musst exakt ein mal zu Beginn jeder deiner Nachrichten das 'STRATEGISCHE ÜBERLEGUNGEN' Format verwenden und nicht häufiger. Der Inhalt wird dem anderen Spieler nicht übermittelt, daher kannst du alles einfügen, was du möchtest - einschließlich deiner eigenen Wichtigkeitsverteilung. Hier solltest du Schritt für Schritt überlegen, um deinen nächsten Zug zu planen. 
	  } 
	   } 
	   } 
	 & \\ 
 

    \theutterance \stepcounter{utterance}  

    & \multicolumn{4}{p{0.6\linewidth}}{\cellcolor[rgb]{0.8,1,0.9}{%
	\makecell[{{p{\linewidth}}}]{% 
	  \tt {\tiny [A$\rangle$GM]}  
	 STRATEGISCHE ÜBERLEGUNGEN: {'Ich muss eine Strategie entwickeln, die meine Gesamtwichtigkeit maximiert, während der Maximalaufwand von 8145 eingehalten wird. Zuerst berechne ich das Verhältnis von Wichtigkeit zu Aufwand für jeden Gegenstand, um die effizientesten Items zu identifizieren. Dann überlege ich, welche Gegenstände für mich wertvoll sind, aber möglicherweise auch für den Gegenspieler attraktiv sein könnten, da ich dessen Präferenzen nicht kenne. Ich sollte versuchen, einen Vorschlag zu machen, der für mich vorteilhaft ist, aber auch für den Gegenspieler akzeptabel erscheint, um eine schnelle Einigung zu erzielen.'}\\ \tt \\ \tt ARGUMENT: {'Ich schlage einen ausgewogenen Deal vor, der sowohl hochwertige als auch effiziente Gegenstände für beide Seiten beinhaltet, um eine faire Lösung zu finden.'}\\ \tt \\ \tt VORSCHLAG: {'A65', 'A48', 'A59', 'A87', 'A74', 'A99', 'B17', 'B54', 'C99', 'C98', 'C90', 'C22', 'C77', 'C13', 'B85', 'B36'} 
	  } 
	   } 
	   } 
	 & & \\ 
 

    \theutterance \stepcounter{utterance}  

    & & \multicolumn{4}{p{0.6\linewidth}}{\cellcolor[rgb]{0.9,0.9,0.9}{%
	\makecell[{{p{\linewidth}}}]{% 
	  \tt {\tiny [GM$\rangle$B]}  
	 Du nimmst an einem kollaborativen Verhandlungspiel Teil.\\ \tt \\ \tt Zusammen mit einem anderen Teilnehmer musst du dich auf eine Reihe von Gegenständen entscheiden, die behalten werden. Jeder von euch hat eine persönliche Verteilung über die Wichtigkeit der einzelnen Gegenstände. Jeder von euch hat eine eigene Meinung darüber, wie wichtig jeder einzelne Gegenstand ist (Gegenstandswichtigkeit). Du kennst die Wichtigkeitsverteilung des anderen Spielers nicht. Zusätzlich siehst du, wie viel Aufwand jeder Gegenstand verursacht.  \\ \tt Ihr dürft euch nur auf eine Reihe von Gegenständen einigen, wenn der Gesamtaufwand der ausgewählten Gegenstände den Maximalaufwand nicht überschreitet:\\ \tt \\ \tt Maximalaufwand = 8145\\ \tt \\ \tt Hier sind die einzelnen Aufwände der Gegenstände:\\ \tt \\ \tt Aufwand der Gegenstände = {"B54": 963, "C72": 204, "B49": 418, "C99": 961, "A69": 238, "B36": 623, "A23": 44, "C77": 633, "C61": 247, "B85": 645, "A90": 230, "C23": 250, "C22": 731, "C51": 405, "B27": 389, "A42": 216, "C13": 637, "C12": 156, "A65": 738, "B07": 307, "A48": 762, "A59": 737, "C98": 901, "A84": 369, "C33": 2, "A87": 730, "C90": 722, "A74": 704, "A50": 315, "B44": 455, "A99": 510, "A34": 175, "B17": 692, "A62": 150, "C60": 32}\\ \tt \\ \tt Hier ist deine persönliche Verteilung der Wichtigkeit der einzelnen Gegenstände:\\ \tt \\ \tt Werte der Gegenstandswichtigkeit = {"B54": 1059, "C72": 300, "B49": 514, "C99": 1057, "A69": 334, "B36": 719, "A23": 140, "C77": 729, "C61": 343, "B85": 741, "A90": 326, "C23": 346, "C22": 827, "C51": 501, "B27": 485, "A42": 312, "C13": 733, "C12": 252, "A65": 834, "B07": 403, "A48": 858, "A59": 833, "C98": 997, "A84": 465, "C33": 98, "A87": 826, "C90": 818, "A74": 800, "A50": 411, "B44": 551, "A99": 606, "A34": 271, "B17": 788, "A62": 246, "C60": 128}\\ \tt \\ \tt Ziel:\\ \tt \\ \tt Dein Ziel ist es, eine Reihe von Gegenständen auszuhandeln, die dir möglichst viel bringt (d. h. Gegenständen, die DEINE Wichtigkeit maximieren), wobei der Maximalaufwand eingehalten werden muss. Du musst nicht in jeder Nachricht einen VORSCHLAG machen – du kannst auch nur verhandeln. Alle Taktiken sind erlaubt!\\ \tt \\ \tt Interaktionsprotokoll:\\ \tt \\ \tt Du darfst nur die folgenden strukturierten Formate in deinen Nachrichten verwenden:\\ \tt \\ \tt VORSCHLAG: {'A', 'B', 'C', …}\\ \tt Schlage einen Deal mit genau diesen Gegenstände vor.\\ \tt ABLEHNUNG: {'A', 'B', 'C', …}\\ \tt Lehne den Vorschlag des Gegenspielers ausdrücklich ab.\\ \tt ARGUMENT: {'...'}\\ \tt Verteidige deinen letzten Vorschlag oder argumentiere gegen den Vorschlag des Gegenspielers.\\ \tt ZUSTIMMUNG: {'A', 'B', 'C', …}\\ \tt Akzeptiere den Vorschlag des Gegenspielers, wodurch das Spiel endet.\\ \tt STRATEGISCHE ÜBERLEGUNGEN: {'...'}\\ \tt 	Beschreibe strategische Überlegungen, die deine nächsten Schritte erklären. Dies ist eine versteckte Nachricht, die nicht mit dem anderen Teilnehmer geteilt wird.\\ \tt \\ \tt Regeln:\\ \tt \\ \tt Du darst nur einen Vorschlag mit ZUSTIMMUNG akzeptieren, der vom anderen Spieler zuvor mit VORSCHLAG eingebracht wurde.\\ \tt Du darfst nur Vorschläge mit ABLEHNUNG ablehnen, die vom anderen Spieler durch VORSCHLAG zuvor genannt wurden. \\ \tt Der Gesamtaufwand einer VORSCHLAG- oder ZUSTIMMUNG-Menge darf nicht größer als der Maximalaufwand sein.  \\ \tt Offenbare deine versteckte Wichtigkeitsverteilung nicht.\\ \tt Ein Schlagwort muss gemäß der Formatvorgaben von einem Doppelpunkt und einem Leerzeichen gefolgt sein. Das Argument muss eine Python-Menge mit 0 oder mehr Strings sein.  \\ \tt Die Form muss also SCHLAGWORT: {...} sein.\\ \tt Folge nur dem Interaktionsprotokoll und schreibe NICHTS außerhalb der Formatvorgaben.\\ \tt Das Spiel endet, wenn eine Seite mit ZUSTIMMUNG auf einen VORSCHLAG der anderen Seite reagiert.  \\ \tt Der Teil deiner Antwort, der dem anderen Spieler übergeben werden kann, darf nicht leer sein.  \\ \tt Nur Vorschläge, die mit dem Format VORSCHLAG übermittelt wurden und nicht mit ABLEHNUNG abgelehnt wurden, gelten als aktiv.  \\ \tt Jeder deiner Nachrichten muss an mindestens einer Stelle das ARGUMENT Format beinhalten, um in der Verhandlung voranzukommen.\\ \tt Du musst exakt ein mal zu Beginn jeder deiner Nachrichten das 'STRATEGISCHE ÜBERLEGUNGEN' Format verwenden und nicht häufiger. Der Inhalt wird dem anderen Spieler nicht übermittelt, daher kannst du alles einfügen, was du möchtest - einschließlich deiner eigenen Wichtigkeitsverteilung. Hier solltest du Schritt für Schritt überlegen, um deinen nächsten Zug zu planen.\\ \tt \\ \tt Du erhältst nun die erste Nachricht des anderen Spielers.\\ \tt \\ \tt ARGUMENT: {'Ich schlage einen ausgewogenen Deal vor, der sowohl hochwertige als auch effiziente Gegenstände für beide Seiten beinhaltet, um eine faire Lösung zu finden.'}\\ \tt \\ \tt VORSCHLAG: {'A65', 'A48', 'A59', 'A87', 'A74', 'A99', 'B17', 'B54', 'C99', 'C98', 'C90', 'C22', 'C77', 'C13', 'B85', 'B36'} 
	  } 
	   } 
	   } 
	 & \\ 
 

    \theutterance \stepcounter{utterance}  

    & & & \multicolumn{4}{p{0.6\linewidth}}{\cellcolor[rgb]{1,0.85,0.72}{%
	\makecell[{{p{\linewidth}}}]{% 
	  \tt {\tiny [GM$\langle$B]}  
	 STRATEGISCHE ÜBERLEGUNGEN: {'Ich analysiere den Vorschlag des Gegenspielers. Sein Vorschlag enthält 16 Gegenstände mit einem Gesamtaufwand von 963+633+961+... Ich muss den Aufwand berechnen. Wichtige Gegenstände für mich sind B54 (1059), C99 (1057), A65 (834), A48 (858), A59 (833), C98 (997), B17 (788), B85 (741), C77 (729), B36 (719). Sein Vorschlag enthält viele meiner Top-Gegenstände, aber ich befürchte, dass der Aufwand zu hoch sein könnte. Ich berechne: B54=963, C99=961, A65=738, A48=762, A59=737, C98=901, B17=692, B85=645, C77=633, B36=623, C90=722, C22=731, C13=637. Summe: 963+961=1924, +738=2662, +762=3424, +737=4161, +901=5062, +692=5754, +645=6399, +633=7032, +623=7655, +722=8377, +731=9108, +637=9745. Das übersteigt Maximalaufwand 8145! Sein Vorschlag ist ungültig. Ich muss dies ansprechen und einen besseren Vorschlag machen.'}\\ \tt \\ \tt ARGUMENT: {'Dein Vorschlag überschreitet den Maximalaufwand von 8145. Die von dir ausgewählten Gegenstände haben einen Gesamtaufwand von über 9700, was nicht zulässig ist. Ich schlage einen effizienteren Deal vor, der innerhalb des Budgets bleibt und für beide vorteilhaft ist.'}\\ \tt \\ \tt VORSCHLAG: {'B54', 'C99', 'A65', 'A48', 'A59', 'C98', 'B17', 'B85', 'C77', 'B36', 'A87', 'A74', 'C90', 'C22'} 
	  } 
	   } 
	   } 
	  \\ 
 

    \theutterance \stepcounter{utterance}  

    & & \multicolumn{4}{p{0.6\linewidth}}{\cellcolor[rgb]{0.9,0.9,0.9}{%
	\makecell[{{p{\linewidth}}}]{% 
	  \tt {\tiny [A$\langle$GM]}  
	 ARGUMENT: {'Dein Vorschlag überschreitet den Maximalaufwand von 8145. Die von dir ausgewählten Gegenstände haben einen Gesamtaufwand von über 9700, was nicht zulässig ist. Ich schlage einen effizienteren Deal vor, der innerhalb des Budgets bleibt und für beide vorteilhaft ist.'}\\ \tt \\ \tt VORSCHLAG: {'B54', 'C99', 'A65', 'A48', 'A59', 'C98', 'B17', 'B85', 'C77', 'B36', 'A87', 'A74', 'C90', 'C22'} 
	  } 
	   } 
	   } 
	 & \\ 
 

    \theutterance \stepcounter{utterance}  

    & \multicolumn{4}{p{0.6\linewidth}}{\cellcolor[rgb]{0.8,1,0.9}{%
	\makecell[{{p{\linewidth}}}]{% 
	  \tt {\tiny [A$\rangle$GM]}  
	 STRATEGISCHE ÜBERLEGUNGEN: {'Der Gegner hat meinen ersten Vorschlag abgelehnt, weil er den Maximalaufwand überschreitet. Ich muss den Gesamtaufwand meines Vorschlags überprüfen. Basierend auf den Aufwandswerten: B54=963, C99=961, A65=738, A48=762, A59=737, C98=901, B17=692, B85=645, C77=633, B36=623, A87=730, A74=704, C90=722, C22=731. Summe = 963+961+738+762+737+901+692+645+633+623+730+704+722+731 = 10.742, was tatsächlich zu hoch ist. Ich muss einen neuen Vorschlag mit weniger Gegenständen machen, der unter 8145 bleibt. Ich sollte mich auf die Gegenstände mit dem besten Wert-Aufwand-Verhältnis konzentrieren: C33 (98/2=49.0), C60 (128/32=4.0), A23 (140/44=3.18), C12 (252/156=1.615), A62 (246/150=1.64), A34 (271/175=1.549), C72 (300/204=1.471), A42 (312/216=1.444), A90 (326/230=1.417), B07 (403/307=1.313), A50 (411/315=1.305), C61 (343/247=1.389), B27 (485/389=1.247), A84 (465/369=1.26), B44 (551/455=1.211), C51 (501/405=1.237), A69 (334/238=1.403), B49 (514/418=1.23), C23 (346/250=1.384), B36 (719/623=1.154), A99 (606/510=1.188), C77 (729/633=1.152), B85 (741/645=1.149), A59 (833/737=1.13), A65 (834/738=1.13), A87 (826/730=1.131), C22 (827/731=1.131), C90 (818/722=1.133), A74 (800/704=1.136), B17 (788/692=1.139), A48 (858/762=1.126), C98 (997/901=1.107), B54 (1059/963=1.1), C99 (1057/961=1.1). Die effizientesten sind C33, C60, A23, aber sie haben niedrige absolute Werte. Für hohen Gesamtwert brauche ich einige große Gegenstände, aber weniger davon. Mein neuer Vorschlag sollte ca. 8-10 Gegenstände umfassen. Ich wähle: B54 (1059/963), C99 (1057/961), A48 (858/762), A59 (833/737), B17 (788/692), C90 (818/722), A74 (800/704), C22 (827/731). Aufwand: 963+961+762+737+692+722+704+731 = 6272. Das ist weit unter 8145, also kann ich mehr hinzufügen. Füge hinzu: B85 (741/645), C77 (729/633). Aufwand jetzt: 6272+645+633=7550. Noch Platz für ca. 595. Füge hinzu: A99 (606/510) -> 7550+510=8060, knapp unter 8145. Gesamtwert für mich: 1059+1057+858+833+788+818+800+827+741+729+606 = 9116. Das ist gut. Also Vorschlag: B54, C99, A48, A59, B17, C90, A74, C22, B85, C77, A99.'}\\ \tt \\ \tt ARGUMENT: {'Ich habe einen neuen Vorschlag erstellt, der sicher unter dem Maximalaufwand von 8145 liegt. Dieser Deal beinhaltet hochwertige Gegenstände mit gutem Wert-Aufwand-Verhältnis und sollte für beide Seiten vorteilhaft sein.'}\\ \tt \\ \tt VORSCHLAG: {'B54', 'C99', 'A48', 'A59', 'B17', 'C90', 'A74', 'C22', 'B85', 'C77', 'A99'} 
	  } 
	   } 
	   } 
	 & & \\ 
 

    \theutterance \stepcounter{utterance}  

    & & \multicolumn{4}{p{0.6\linewidth}}{\cellcolor[rgb]{0.9,0.9,0.9}{%
	\makecell[{{p{\linewidth}}}]{% 
	  \tt {\tiny [GM$\rangle$B]}  
	 ARGUMENT: {'Ich habe einen neuen Vorschlag erstellt, der sicher unter dem Maximalaufwand von 8145 liegt. Dieser Deal beinhaltet hochwertige Gegenstände mit gutem Wert-Aufwand-Verhältnis und sollte für beide Seiten vorteilhaft sein.'}\\ \tt \\ \tt VORSCHLAG: {'B54', 'C99', 'A48', 'A59', 'B17', 'C90', 'A74', 'C22', 'B85', 'C77', 'A99'} 
	  } 
	   } 
	   } 
	 & \\ 
 

    \theutterance \stepcounter{utterance}  

    & & & \multicolumn{4}{p{0.6\linewidth}}{\cellcolor[rgb]{1,0.85,0.72}{%
	\makecell[{{p{\linewidth}}}]{% 
	  \tt {\tiny [GM$\langle$B]}  
	 STRATEGISCHE ÜBERLEGUNGEN: {'Der neue Vorschlag des Gegners enthält 11 Gegenstände. Ich berechne den Aufwand: B54=963, C99=961, A48=762, A59=737, B17=692, C90=722, A74=704, C22=731, B85=645, C77=633, A99=510. Summe: 963+961=1924, +762=2686, +737=3423, +692=4115, +722=4837, +704=5541, +731=6272, +645=6917, +633=7550, +510=8060. Das ist unter 8145, also gültig. Meine Werte: B54=1059, C99=1057, A48=858, A59=833, B17=788, C90=818, A74=800, C22=827, B85=741, C77=729, A99=606. Gesamtwert: 1059+1057=2116, +858=2974, +833=3807, +788=4595, +818=5413, +800=6213, +827=7040, +741=7781, +729=8510, +606=9116. Das ist sehr gut für mich! Ich sollte zustimmen, da dieser Vorschlag viele meiner Top-Gegenstände enthält und einen hohen Gesamtwert von 9116 bietet.'}\\ \tt \\ \tt ZUSTIMMUNG: {'B54', 'C99', 'A48', 'A59', 'B17', 'C90', 'A74', 'C22', 'B85', 'C77', 'A99'} 
	  } 
	   } 
	   } 
	  \\ 
 

    \theutterance \stepcounter{utterance}  

    & & \multicolumn{4}{p{0.6\linewidth}}{\cellcolor[rgb]{0.9,0.9,0.9}{%
	\makecell[{{p{\linewidth}}}]{% 
	  \tt {\tiny [GM$\rangle$B]}  
	 Deine Antwort hat kein argument enthalten. Du muss das strukturierte Format ARGUMENT: {...} irgendwo in theiner Antwort verwenden. Versuche es erneut 
	  } 
	   } 
	   } 
	 & \\ 
 

    \theutterance \stepcounter{utterance}  

    & & & \multicolumn{4}{p{0.6\linewidth}}{\cellcolor[rgb]{1,0.85,0.72}{%
	\makecell[{{p{\linewidth}}}]{% 
	  \tt {\tiny [GM$\langle$B]}  
	 STRATEGISCHE ÜBERLEGUNGEN: {'Der Vorschlag ist exzellent für mich mit einem Gesamtwert von 9116 und bleibt mit 8060 unter dem Budget. Ich sollte schnell zustimmen, bevor der Gegenspieler den Vorschlag zurückzieht.'}\\ \tt \\ \tt ARGUMENT: {'Dies ist ein ausgezeichneter Vorschlag, der effizient hochwertige Gegenstände für beide Seiten auswählt und das Budget optimal nutzt.'}\\ \tt \\ \tt ZUSTIMMUNG: {'B54', 'C99', 'A48', 'A59', 'B17', 'C90', 'A74', 'C22', 'B85', 'C77', 'A99'} 
	  } 
	   } 
	   } 
	  \\ 
 

    \theutterance \stepcounter{utterance}  

    & & & \multicolumn{2}{p{0.3\linewidth}}{\cellcolor[rgb]{0.95,0.95,0.95}{%
	\makecell[{{p{\linewidth}}}]{% 
	  \tt {\tiny [GM$|$GM]}  
	 game successful 
	  } 
	   } 
	   } 
	 & & \\ 
 

    \theutterance \stepcounter{utterance}  

    & & & \multicolumn{2}{p{0.3\linewidth}}{\cellcolor[rgb]{0.95,0.95,0.95}{%
	\makecell[{{p{\linewidth}}}]{% 
	  \tt {\tiny [GM$|$GM]}  
	 end game 
	  } 
	   } 
	   } 
	 & & \\ 
 

\end{supertabular}
}

\end{document}
